\section{Degree-one action and the quadratic analytic Aeppli obstruction}
\label{sec:r23-quadratic-aeppli-action}

\paragraph{Status and exact gain.}
This section is a \texttt{DIRECT\_COROLLARY} of the frozen R6
degree-one action and piecewise-smooth Leibniz theorem.  It is not
counted as a new principal theorem.  It makes one distinction in
the R22 singular interface precise: in formal degree two, a full
\emph{intrinsic} Aeppli identity already suffices for annihilation
by every allowed original background.  There is no additional
arbitrary-background action hypothesis for that implication.
The existence of the identity remains \texttt{OPEN}.  The analogous
full-polar compatibility and the unchanged KC12 seminorm conclusion
are not proved here.

Fix the original finite extension $K/\mathbf Q_p$ and its original
absolute value.  Let $U$ be an original good, Hausdorff, boundaryless,
quasi-smooth analytic space, locally of pure dimension $n$.  All
sheaves, restrictions and germs below use its ordinary topology.
Write $B^{a,b}$ for the genuine analytic Mihatsch delta-form sheaf,
including all actual moment refinements.  Use standard CLD signs:
$d'd''=-d''d'$ and $D=d'd''$.  Let $S$ be any original symmetric
middle-positive, separately $d'$- and $d''$-closed current of degree
$(p,p)$, with $0\le p\le n-2$.  No conormality, model origin,
weak-approximation, or special decomposition is assumed for $S$.

\subsection{The precise first-degree multiplication input}

The frozen R6 theorem
\texttt{thm:r6-arbitrary-background-degree-one-delta} supplies a
canonical linear operation, denoted here by
\[
 M_S:B^{1,1}(U)\longrightarrow\mathcal D'^{p+1,p+1}(U),
 \qquad M_S(\alpha)=\alpha\wedge S.                 \tag{QA1}
\]
Its values are coefficient-Radon currents.  It respects ordinary
restriction, gluing, every actual refinement, and actual equality
in $B^{1,1}$.  It agrees with the already established multiplication
by piecewise-smooth forms and satisfies
\[
       M_S(fD_\delta g)=fD(gS)                       \tag{QA2}
\]
for continuous piecewise-smooth scalars $f,g$ in their finite local
presentations; $g$ is expanded into the permitted local continuous
psh differences.  The required local decomposition is the frozen
\texttt{lem:r6-continuous-pws-dsh}.  Only this degree-one operation
is used.  In particular, no expression $(D\alpha)\wedge S$ is
defined by appealing to a presumed higher-degree action.

Every actual form of total bidegree $(0,1)$ is piecewise smooth,
and locally has an anchored finite presentation
\[
                    \beta=\sum_j f_jd''_Pg_j,        \tag{QA3}
\]
with continuous piecewise-smooth scalars $f_j,g_j$.  This is the
codimension-zero module generation used explicitly in the proof of
R6, \texttt{cor:r6-all-background-qone-full-delta}.  A form of
bidegree $(1,0)$ has the column-exchanged presentation.
Their products with $S$ are the frozen canonical PS products.

\begin{lemma}[Full first derivatives commute with the available action]
For every actual $\beta\in B^{0,1}(U)$ and
$\gamma\in B^{1,0}(U)$,
\[
 \begin{split}
 M_S(d'\beta)&=d'(\beta\wedge S),\\
 M_S(d''\gamma)&=d''(\gamma\wedge S).
 \end{split}                                          \tag{QA4}
\]
These are equalities of ordinary currents on all original smooth
compact tests, with all full polyhedral boundary contributions
retained.  They hold in every degree $p<n$ where (QA1) is defined;
the stronger restriction $p\le n-2$ is needed only for the
subsequent nontrivial quadratic derivative.
\end{lemma}

\begin{proof}
The assertion is local in the ordinary topology, so use (QA3).
For one summand the actual full delta Leibniz rule reads
\[
 d'(f d''_P g)=d'_P f\wedge d''_P g+fD_\delta g.
                                                        \tag{QA5}
\]
The scalar has no boundary differential, and $d''g=d''_Pg$;
the singular derivative of the latter is exactly $D_\delta g$.
Applying (QA1)--(QA2), including agreement with PS multiplication,
gives
\[
 M_S\bigl(d'(f d''_P g)\bigr)
       =d'_P f\wedge d''_P g\wedge S+fD(gS).             \tag{QA6}
\]

The frozen first-derivative identity is
$d''(gS)=d''_Pg\wedge S$.  The right side is in the tangential
coefficient-Radon module of R6.  Its $d'$ derivative is $D(gS)$,
which is coefficient-Radon because $g$ is a finite local continuous
psh difference.  Therefore the frozen
\texttt{lem:r6-pws-radon-leibniz}, applied to multiplication by
$f$, gives
\[
 d'\bigl(f(d''_P g\wedge S)\bigr)
       =d'_P f\wedge d''_P g\wedge S+fD(gS).             \tag{QA7}
\]
Equations (QA6)--(QA7) agree.  Sum the finitely many terms and
use ordinary locality to obtain the first equation of (QA4).
The identical argument with the two columns exchanged gives the
second equation, with the standard CLD anticommutation signs.

For clarity, the proof in R6 subsequently integrates (QA7) when
$p=n-1$.  That degree assumption is used only for the top-degree
integral there, not for the local current identity (QA7).  No
integration of a non-top-degree form is performed here.
The input Leibniz theorem proves equality on the complete original
test sheaf by convergence against fixed Radon coefficient measures;
we do not delete a boundary term on a charged face.  Since all
products on both sides are already intrinsic, the equations respect
every actual refinement, including deficient-rank moment maps.
\end{proof}

\subsection{An intrinsic obstruction map for an arbitrary background}

Let $\mathcal E^{1,1}$ be the image sheaf
\[
 \mathcal E^{1,1}
   =d'B^{0,1}+d''B^{1,0}\ \subset B^{1,1}.               \tag{QA8}
\]
Here an image sheaf means that the decomposition is required locally
in the ordinary topology; a global decomposition on $U$ is not
assumed.  Nilpotence and anticommutation give $D\mathcal E^{1,1}=0$.
Define the actual degree-$(1,1)$ Aeppli sheaf by
\[
 \mathcal H_A^{1,1}(B)
   =\ker(D:B^{1,1}\to B^{2,2})/\mathcal E^{1,1}.          \tag{QA9}
\]
This definition uses the original analytic delta quotient itself.
It is not the analogous quotient on one chosen ambient real space,
nor a quotient defined only by the tests extending to one marked
point.  It introduces no new carrier seminorm or relation generator.

\begin{proposition}[The available action factors through the actual quotient]
The operation
\[
        \mathfrak o_S(\alpha):=D\bigl(M_S(\alpha)\bigr)   \tag{QA10}
\]
annihilates $\mathcal E^{1,1}$.  It therefore induces a linear
map of ordinary sheaves
\[
 \mathfrak o_S:\mathcal H_A^{1,1}(B)
                  \longrightarrow\mathcal D'^{p+2,p+2}. \tag{QA11}
\]
Its values are separately closed currents.  Formula (QA10) has
support contained in $\operatorname{supp}\alpha\cap
\operatorname{supp}S$ for every representative $\alpha$.
All original tests and actual refinements are retained.
No assertion that (QA11) is zero on its entire domain is made.
\end{proposition}

\begin{proof}
By (QA4), the value of (QA10) on $d'\beta$ is
$d'd''d'(\beta\wedge S)=0$.  Its value on $d''\gamma$ is
$d'd''d''(\gamma\wedge S)=0$.  These identities hold for
distributions with the indicated standard signs.  Ordinary locality
annihilates the image sheaf in (QA8), so the map is well defined
even when representatives differ only by local exact expressions.
Restriction of (QA10) to the kernel in (QA9) proves (QA11).
In fact (QA10) already factors through $B^{1,1}/\mathcal E^{1,1}$;
the restriction is made to isolate the actual zero-relation problem.

Both derivatives of $D M_S(\alpha)$ vanish by nilpotence and
anticommutation.  The support assertion follows from the support
theorem for (QA1) and the support property of a distributional
derivative.  The operation $D M_S(\alpha)$ need not be
coefficient-Radon for arbitrary $\alpha$; no such general claim is
needed.  For the finite Hessian-product expressions below it is the
already defined finite signed BT current, hence coefficient-Radon.
Every definition used is a morphism in the actual original sheaves,
so their all-test and all-refinement compatibilities pass to the
quotient.  An ordinary-current-to-GK-test extension is not inferred.
\end{proof}

\subsection{The quadratic curvature transgression}

Let $\mathscr P$ be the real sheaf of finite continuous PL functions
in actual original unit moments, allowing finite sums and ordinary
restriction.  These functions belong locally to the frozen DC
class.  For the formal symmetric square, over $\mathbf R$ and not
over the ring of pointwise scalar functions, define
\[
 \begin{split}
 \widetilde\tau([h][g])
       &=\tfrac12\bigl(hD_\delta g+gD_\delta h\bigr)
                                                    \in B^{1,1},\\
 \kappa_2([h][g])&=D_\delta h\wedge D_\delta g
                                                    \in B^{2,2}.
 \end{split}                                           \tag{QA12}
\]
These are linear maps from $\operatorname{Sym}^2_{\mathbf R}
\mathscr P$.  Full delta Leibniz and the separate closedness of a
scalar curvature give
\[
                         D\widetilde\tau=\kappa_2.      \tag{QA13}
\]
Thus every actual analytic quadratic relation $P\in\ker\kappa_2$
has an intrinsic class
\[
                   \tau(P)=[\widetilde\tau(P)]
                                   \in\mathcal H_A^{1,1}(B).
                                                        \tag{QA14}
\]
One may instead use an ordered primitive $\alpha_P=\sum_\nu
a_\nu h_\nu D_\delta g_\nu$ for a displayed polynomial
$P=\sum_\nu a_\nu[h_\nu][g_\nu]$.  The class is unchanged,
because the exact full identity
\[
 hD_\delta g-gD_\delta h
       =d'(h d''g)+d''(g d'h)                           \tag{QA15}
\]
has the indicated signs.  Indeed the two first-derivative cross
terms cancel, and $d''d'h=-D_\delta h$.

\begin{theorem}[Quadratic intrinsic exactness suffices for every allowed action]
Let $P=\sum_\nu a_\nu[h_\nu][g_\nu]$ be a finite original
PL presentation on one ordinary representative of a germ, and
assume $\kappa_2(P)=0$ in its actual analytic delta sheaf.
For every allowed original $S$ of degree $p\le n-2$,
\[
          \mathcal A_S(P)=\mathfrak o_S\bigl(\tau(P)\bigr).
                                                        \tag{QA16}
\]
In particular, if on that ordinary germ
\[
           \alpha_P=d'\beta+d''\gamma,
       \quad \beta\in B^{0,1}=PS^{0,1},\quad
             \gamma\in B^{1,0}=PS^{1,0},                \tag{QA17}
\]
then $\mathcal A_S(P)=0$ for every allowed original $S$, on
every original compact smooth complementary test and every actual
refinement.  The absolute carrier of $S$ is unchanged.
The ordinary representative for (QA17) can be chosen before
the background and the tests, since it is an intrinsic identity
of the specified relation.
\end{theorem}

\begin{proof}
Expand each signed PL potential into its fixed finite convex-PL
differences.  R6 and (QA2) identify
\[
          M_S(\alpha_P)=\sum_\nu a_\nu h_\nu D(g_\nu S).
                                                        \tag{QA18}
\]
Applying $D$ gives exactly the recursively defined relative
quadratic BT action.  The frozen signed multilinearity and symmetry
identify the same expression for the symmetrized primitive in
(QA12).  Alternatively, (QA15) and (QA4) show directly that the
two values after $D$ agree.  This proves (QA16).

Under (QA17), (QA4) gives
\[
 M_S(\alpha_P)=d'(\beta\wedge S)+d''(\gamma\wedge S).
                                                        \tag{QA19}
\]
Apply $D$ and use nilpotence to obtain zero.  This reasoning takes
place on the same ordinary representative as the identity and
uses only its already intrinsic degree-one products.  Arbitrary
compact original tests in that representative are allowed, including
tests whose own units do not extend to the germ center.
All actual refinements were part of the input identities.
The support statements for (QA1) and ordinary derivatives show
that no replacement or enlargement of the original carrier has
entered the argument.  If the identity is only germwise, choose
one representative where it holds before choosing $S$ or a test;
no global polar is transported to that smaller quotient.
\end{proof}

The finite-tail version is immediate in its exact limited scope.
If one fixed common tuple of continuous psh potentials is appended
to every term of $P$, first form its frozen positive separately
closed BT product with $S$ and apply the theorem to that background.
Fixed signed finite PL tails are treated by their finite positive
expansions.  Distinct termwise tail tuples are not replaced by one
common background, and this observation does not settle the general
higher-degree coupled relation.

\subsection{What remains open and what this implication does not supply}

The new readback removes an additional action hypothesis only from
the \emph{quadratic arbitrary-background} route through (QA17).
The actual existence question is now explicit: for the restricted
class $\alpha_P$ arising from an actual quadratic PL-zero, must
$\tau(P)$ vanish in the actual Aeppli sheaf, or at least be killed
by every map (QA11)?  This is \texttt{OPEN}.  The frozen CL
calculation replaces $\alpha_P$, modulo full exact types, by the
specified finite piecewise-constant PS energy form; PH's pointed
homotopy retains a singular boundary and does not prove its
exactness.  Actual $D$-closedness alone is not inserted as a new
kernel of (QA10).

A general full original carrier polar is not an allowed positive
current $S$.  Consequently neither (QA1) nor (QA4) can be applied
to it by a change of notation.  Even if (QA17) were proved for
all quadratic relations, the theorem above would establish their
arbitrary-background annihilation, not automatically zero of the
unchanged carrier-CLN seminorm.  The frozen IF strict gap prohibits
inferring that seminorm assertion solely from carried-current
annihilation.  Conversely, the theorem supplies no realization of
an arbitrary real conormal current as a full original polar.

The differentials in (QA4), (QA15), and (QA17) are the full analytic
delta differentials.  Replacing them by $d'_P,d''_P$ would delete
the singular curvature in (QA5) and invalidate (QA19).  No
convolution of an actual zero is asserted to remain zero.  No
ordinary-current equality is promoted to equality on all GK delta
tests.  The relevant degree-one delta operation is an explicit
frozen theorem, and no degree-two delta operation is assumed in
order to prove its own presentation independence.  Thus KC12,
OPEN-GQ in general, and variable-PS/CTO remain \texttt{OPEN}.

\paragraph{Dependencies and participation.}
The root constructed this quotient/action deduction after directly
reading the frozen R6 degree-one action, its anchored generation,
and its Radon Leibniz proof.  The audit researcher independently
checked the typing of those frozen inputs and confirmed that the
local product identity precedes the special top-degree integral;
that input calibration is disclosed as a contribution, not counted
as an independent review of this source.  The exact source and input
hashes, any subsequent author-separated review, and the R22-to-R23
scope correction are recorded separately.  Internal review is not
external peer review or formal verification.
